\section{材料与数据}\label{ch:data}

\subsection{研究区概况}

【填写研究区位置、自然条件、实验对象和数据采集背景。若论文不是地理或农业方向，可将本节改为“实验对象与样本来源”。】

研究区位于【省/市/县或实验平台名称】，地理范围为【经纬度或空间范围】。研究对象为【作物/材料/系统/设备】，数据采集时间为【日期范围】。研究区或实验对象的基本情况见表\ref{tab:study-area}。

\begin{table}[ht]
  \centering
  \caption{研究区或实验对象基本信息}
  \label{tab:study-area}
  \begin{tabular}{lll}
    \toprule
    项目 & 内容 & 说明 \\
    \midrule
    研究对象 & 【填写】 & 【如作物类型、样本类型】 \\
    时间范围 & 【填写】 & 【如 20XX 年 X 月至 X 月】 \\
    数据规模 & 【填写】 & 【如样本数、影像数、观测次数】 \\
    空间尺度 & 【填写】 & 【如像元、样方、实验单元】 \\
    \bottomrule
  \end{tabular}
\end{table}

\subsection{数据来源}

【填写原始数据来源、仪器设备、公开数据集或人工标注流程。需要说明数据格式、分辨率、采样频率、坐标系统或实验条件。】

本文使用的数据包括【数据 A】、【数据 B】和【辅助数据】。其中，【数据 A】用于表征【输入信息】，【数据 B】用于提供【标签或验证信息】，【辅助数据】用于【质量控制或空间配准】。

\begin{longtable}{lll}
\caption{主要数据源及用途}\label{tab:data-source}\\
\toprule
数据类型 & 主要内容 & 用途 \\
\midrule
\endfirsthead
\toprule
数据类型 & 主要内容 & 用途 \\
\midrule
\endhead
\bottomrule
\endfoot
【数据 A】 & 【分辨率/字段/格式】 & 【特征提取】 \\
【数据 B】 & 【分辨率/字段/格式】 & 【标签构建或验证】 \\
【辅助数据】 & 【分辨率/字段/格式】 & 【筛选、配准或解释】 \\
\end{longtable}

\subsection{数据预处理}

【填写预处理流程。建议按实际执行顺序写：格式转换、异常值剔除、坐标统一、空间/时间匹配、归一化、样本筛选。】

预处理包括以下步骤：首先，对原始数据进行格式统一和异常值检查；其次，将不同来源数据配准到统一坐标或样本索引；最后，根据质量控制规则剔除无效样本，形成可用于建模的样本表。

设原始观测值为 $x$，归一化后的特征为 $x'$，本文采用的 min--max 归一化公式为
\begin{equation}
x' = \frac{x - x_{\min}}{x_{\max} - x_{\min}}.
\end{equation}

其中，$x_{\min}$ 和 $x_{\max}$ 分别表示训练集中该特征的最小值和最大值。验证集和测试集仅使用训练集统计量进行变换，以避免数据泄露。
